\section*{الفصل \ref{ch:b3:rna-regulation} --- RNA غير المشفِّر والتنظيم بعد الاستنساخ}

\begin{solution}{exo:b3:rna-regulation:1}
miRNA: مشفَّر في الجينوم دبوسَ شعر، طوله \qty{22}{nt}، يقطعه \omterm{def:b3:rna-regulation:mirna}{Drosha}
ثم \omterm{def:b3:rna-regulation:rnai}{Dicer}، ويزاوج مزاوجة ناقصة ببذرته فيكبت ترجمة رسائل كثيرة
وثباتها. \omterm{def:b3:rna-regulation:ncrna}{وsiRNA}: من RNA طويل مزدوج الطاق (فيروسي أو نقيل أو تجريبي)،
طوله \qty{21}{nt}، يعتمد على \omterm{def:b3:rna-regulation:rnai}{Dicer}، ويزاوج مزاوجة تامة فيوجّه
\omterm{def:b3:rna-regulation:rnai}{Argonaute} إلى قطع الهدف. \omterm{def:b3:rna-regulation:ncrna}{وpiRNA}: من نسخ مفردة الطاق \omterm{def:b3:rna-regulation:pirna}{لعناقيد piRNA}،
طوله \qtyrange{26}{31}{nt}، لا يعتمد على \omterm{def:b3:rna-regulation:rnai}{Dicer}، ويُحمَّل على بروتينات
PIWI في السلالة الجنسية، ويقطع نسخ العناصر القافزة (البنغ بونغ)
ويوجّه إسكات \omterm{def:b3:chromatin-epigenetics:nucleosome}{الصبغين} على DNA تلك العناصر.
\end{solution}

\begin{solution}{exo:b3:rna-regulation:2}
يستنسخ بوليميراز II جزيء pri-miRNA؛ ويقتطع \omterm{def:b3:rna-regulation:mirna}{Drosha} (مع DGCR8) الدبوسَ
في النواة؛ ويصدّر exportin-5 جزيء pre-miRNA؛ ويقطع \omterm{def:b3:rna-regulation:rnai}{Dicer} العروة
فتبقى ثنائية طولها \qty{22}{nt}؛ ويُحمَّل أحد الطاقين في
\omterm{def:b3:rna-regulation:rnai}{Argonaute}؛ وتزاوج البذرة (النوكليوتيدات 2--8) موقعًا في UTR 3$'$؛
ويستدعي \omterm{def:b3:rna-regulation:rnai}{Argonaute} آلات نزع الأدنين ونزع القلنسوة ويحجب البدء ---
فتُترجم الرسالة أقل وتتحلل أسرع.
\end{solution}

\begin{solution}{exo:b3:rna-regulation:3}
كان RNA المصنوع خارج الجسم ببوليميراز عاثية يحوي ملوثات مزدوجة الطاق
(فالبوليميراز ينسخ الطاق الآخر من القالب ويتجاوز النهايات أيضًا)،
فحملت مستحضرات ``ذات المعنى'' المحرّضَ الحقيقي. ونقّى فاير وميلو كل
طاق، وبيّنا أن كلًا منهما وحده يفعل القليل، ثم لحّماهما عمدًا: فكان
RNA المزدوج أقوى مئة مرة على الأقل، إذ يكفي منه بضعة جزيئات في
الخلية.
\end{solution}

\begin{solution}{exo:b3:rna-regulation:4}
الخلية المجوَّعة من الحديد: يرتبط IRP بعناصر IRE. الفيريتين: يُحجب
بدء الترجمة فلا يُصنع بروتين الخزن (\omterm{prop:b3:rna-regulation:translation}{تحكم في الترجمة}). ومستقبل
الترانسفيرين: تُحمى رسالته من نوكليازها فتتراكم، فيُصنع مستقبل أكثر
ويُستورد حديد أكثر (تحكم في ثبات الرسالة).
\end{solution}

\begin{solution}{exo:b3:rna-regulation:5}
$\gamma = \ln 2/\qty{120}{min} = \qty{0.0058}{min^{-1}}$؛ و$m^{*} =
5/0.0058 \approx 870$ جزيئًا. وبمضاعفة التحلل: $m^{*} \approx 430$،
وزمن النصف $\qty{120}{min}/2 = \qty{60}{min}$.
\end{solution}

\begin{solution}{exo:b3:rna-regulation:6}
مجموع تسلسل UTR 3$'$ هو $2\times 10^{7}$ نوكليوتيد؛ والمطابقات
المصادفة $2\times
10^{7}/4^{7} = 2\times 10^{7}/\num{16384} \approx \num{1200}$ موقع.
ويميز الهدفَ الحقيقي حفظُ الموقع عبر الأنواع، والسياق الملائم (جوار
غني بالأدنين واليوراسيل، وموقع بعيد عن رامزة الوقف، ومزاوجة إضافية عند النهاية
3$'$ من miRNA)، والتعبير المشترك عن miRNA والهدف في الخلية نفسها،
والتجربة: إذ تهبط الرسالة حين يحضر miRNA ولا تهبط حين يُطفر الموقع.
\end{solution}

\begin{solution}{exo:b3:rna-regulation:7}
(أ) XX مع \emph{tra} معطّلة: يُصنع SXL ولا يُصنع TRA، فتأخذ
\emph{dsx} التشذيب الذكري الافتراضي --- ذكر جسدي (عقيم). (ب) XY مع
TRA بنيوي: TRA-2 حاضر في الجنسين، فيُصنع DSX في صورته الأنثوية ---
تطور جسدي أنثوي (``شبه أنثى'' عقيمة، والسلالة الجنسية تتبع إشارات
أخرى). (ج) بلا \emph{dsx}: فلا هذه الصورة ولا تلك، فلا يُفرض أي
برنامج من برامج التمايز النهائي --- كائن خنثوي بصفات مختلطة أو
ناقصة.
\end{solution}

\begin{solution}{exo:b3:rna-regulation:8}
عند الحالة المستقرة تتناسب الرسالة، ومن ثم البروتين، مع عمر نصف
الرسالة: $\qty{180}{min}/\qty{15}{min} = 12$ ضعفًا من البروتين.
والبروتين إشارة نمو سوية يُراد لها أن تكون عابرة، دفقةً بعد كل
منبّه؛ فإذا صُنع باثني عشر ضعفًا وباستمرار، دفع التكاثر بلا منبّه.
فالجرعة والتوقيت، لا التسلسل، هما ما يجعله مسرطنًا.
\end{solution}

\begin{solution}{exo:b3:rna-regulation:9}
(أ) أم مخبرية وأب بري: لا تحمل البيوض \omterm{def:b3:rna-regulation:ncrna}{piRNA} ضد عناصر P (فسلالة الأم
لم تلقها قط)، فتقفز عناصر P الأبوية بحرية في السلالة الجنسية للذرية
--- فتكون عقيمة (خللية التخلّق). (ب) أم برية: تحوي بيوضها \omterm{def:b3:rna-regulation:ncrna}{piRNA} ضد
عناصر P فتُسكتها في الذرية --- فتكون خصبة. وعدم التناظر هو إيداع
الأم \omterm{def:b3:rna-regulation:ncrna}{piRNA}، أي توريث RNA لا توريث مورّثات.
\end{solution}

\begin{solution}{exo:b3:rna-regulation:10}
يتكامل البروتين على رسالته على مدى عمره $\tau_{p}$. والرسائل تعيش
$\tau_{m} = 1/(\gamma + \gamma' R)$؛ فخلال $\tau_{p}$ يرى البروتين
نحو $\tau_{p}/\tau_{m}$ من أعمار الرسائل المستقلة، كل منها حدث ولادة
وموت عشوائي. والتقلب النسبي لمتوسط على $N$ حدثًا مستقلًا يهبط
مثل $1/\sqrt{N}$، فإن $\tau_{m}$ الأقصر تحت تحكم \omterm{def:b3:rna-regulation:ncrna}{RNA المكروي} --- عند
متوسط الرسالة نفسه، وهو يقتضي $k$ أعلى بالنسبة نفسها ---
يعطي بروتينًا أنعم. وتدفع الخلية ثمن \omterm{def:b3:rna-regulation:ncrna}{RNA المكروي} والاستنساخ الزائد
لتشتري الدقة والسرعة: فالمستوى الذي يضبطه توازن بين تركيب سريع
وتحلل سريع أقل ضجيجًا وأسرع في إعادة الضبط من المستوى نفسه إذا
ضبطه تركيب بطيء وتحلل بطيء.
\end{solution}

\begin{solution}{exo:b3:rna-regulation:11}
الدودة: يصنع بوليميراز RNA المعتمد على RNA جزيئاتِ \omterm{def:b3:rna-regulation:ncrna}{siRNA} جديدة من
رسالة الهدف نفسها، فتتجدد الإشارة ما دام الهدف يُستنسخ، وتنجو من
التمديد عبر الانقسامات، ومع قناة تمرّ RNA المزدوج بين الخلايا تنتشر
إلى أنسجة أخرى وإلى السلالة الجنسية بضعة أجيال. والثدييات: لا تضخيم،
فتُستهلك جزيئات \omterm{def:b3:rna-regulation:ncrna}{siRNA} وتُمدَّد عند كل انقسام، ولا تعمل إلا في الخلايا
التي تلقتها، وتخبو في أيام في الخلايا المنقسمة. والنتيجة: أن دواء
\omterm{def:b3:rna-regulation:ncrna}{siRNA} يجب أن يوصَل إلى كل خلية هدف، وأنه أنجع في الخلايا التي لا
تنقسم (الخلايا الكبدية، ومن هنا أدوية الكبد)، وأنه يحتاج إلى جرعات
متكررة.
\end{solution}

\begin{solution}{exo:b3:rna-regulation:12}
(أ) RNA وظيفي: هدم النسخة بعد صنعها (بجزيء \omterm{def:b3:rna-regulation:ncrna}{siRNA} أو قليل نوكليوتيد مضاد
للمعنى) يعطي نمطًا ظاهريًا، والتعبير عن RNA من مورث نقيل في مكان آخر
من الجينوم ينقذه. (ب) استنساخ وظيفي: إدخال إشارة بَلْمَرَة أدنين
باكرة توقف الاستنساخ (أو حذف المروّج) يعطي النمط الظاهري، في حين لا
يعطيه هدمُ RNA بعد الاستنساخ، ولا ينقذ المورث النقيل المنبتّ --- فالمهم
فعل الاستنساخ (أو عنصر DNA)، كما في كثير من RNA المقترنة
بالمعزِّزات. (ج) ضجيج: لا تغيّر أي من المعالجتين شيئًا، وRNA بأقل من
نسخة في الخلية، وتسلسله وتعبيره غير محفوظين.
\end{solution}

\begin{solution}{pb:b3:rna-regulation:1}
\textbf{1.} $\gamma = \ln 2/40 = \qty{0.0173}{min^{-1}}$؛ و$m^{*} =
3/0.0173 \approx 170$ رسالة.
\textbf{2.} $\delta_{p} = \ln 2/180 = \qty{0.00385}{min^{-1}}$؛ و$P^{*} =
\beta m^{*}/\delta_{p} = 2\times 173/0.00385 \approx 9.0\times 10^{4}$
بروتين.
\textbf{3.} التحلل $4\gamma = \qty{0.069}{min^{-1}}$؛ و$m^{*} = 3/0.069
\approx 43$؛ وثابت الزمن $1/0.069 = \qty{14}{min}$.
\textbf{4.} عند $\beta = 1$: $P^{*} = 43/0.00385 \approx 1.1\times 10^{4}$؛
فالكبت $8$ أضعاف ($4$ من التحلل و$2$ من الترجمة).
\textbf{5.} زمن نصف الرسالة $\ln 2/(4\gamma) = \qty{10}{min}$. ثم
يتحلل البروتين بعمر نصفه هو، \qty{3}{h}: فتحلل البروتين هو ما يحدّ
المفتاح.
\textbf{6.} نعم: تهبط الرسالة إلى $1/4$ فقط (وهي ما تزال سهلة
الكشف) بينما يهبط البروتين إلى $1/8$ ويستمر في الهبوط مع دوران
البروتين القديم؛ والمشاهدة الأولى بأن البروتين زال ``والرسالة
باقية'' كانت تعكس أساسًا الشطر الترجمي من الكبت.
\textbf{7.} $10^{6}/250 = 4000$ جزيء في البيضة.
\textbf{8.} $4000/550 \approx 7$ في الخلية عند الفقس. و$4000/2^{n} < 1$
عند $n \ge 12$ انقسامًا من اللاقحة --- أي نحو ثلاثة انقسامات
بعد الفقس.
\textbf{9.} ينسخ البوليميراز رسالة الهدف إلى RNA مزدوج جديد، يُقطع
إلى \omterm{def:b3:rna-regulation:ncrna}{siRNA} ثانوية: فيتجدد المحرّض من الهدف نفسه، فينجو من التمديد
وينتشر (إذ تمرّ قناة RNA المزدوج بين الخلايا وإلى السلالة
الجنسية). ويخبو لأن التضخيم يحتاج إلى الهدف وإلى محرّض أولي معًا؛
فمتى زال RNA المحقون مُدّد المخزون الثانوي جيلًا بعد جيل وصُفّرت
السلالة الجنسية.
\textbf{10.} $5000/2^{n} < 100$: أي $2^{n} > 50$، ومنه $n = 6$ أيام ($78$
نسخة باقية).
\textbf{11.} الخلايا التي لا تنقسم لا تمدّد \omterm{def:b3:rna-regulation:rnai}{RISC}؛ والحدّ هو العمر
الكيميائي لجزيء \omterm{def:b3:rna-regulation:ncrna}{siRNA} (تحلل بالنوكليازات، يبطئه تعديل الطاقين
كيميائيًا) والرشح البطيء من المستودع الجسيمي الذي يغذي الهيولى.
\textbf{12.} $2\times 10^{7}/\num{16384} \approx \num{1200}$ مطابقة
مصادفة. والوراثة هي التي حسمت: ففقد \emph{lin-14} يعطي النمط
الظاهري المعاكس لفقد \emph{lin-4}، وأليلات \emph{lin-14} كسبية
الوظيفة كانت حذوفًا لمواقع UTR 3$'$ بعينها، والمواقع السبعة كانت
محفوظة.
\textbf{13.} المستوى $\propto$ عمر النصف: $50/10 = 5$ أضعاف ارتفاعًا.
\textbf{14.} $1/0.05 = 20$ ضعفًا هبوطًا في تركيب الفيريتين.
\textbf{15.} نسبة الاستيراد إلى الخزن ترتفع $5\times 20 = 100$ ضعف من
الحالة الغنية بالحديد إلى الفقيرة به.
\textbf{16.} لا يتجمّع العنقود إلا حين يتوفر الحديد الهيولي،
فيكون انتظام البروتين قراءةً مباشرة لتركيز الحديد، بلا مستقبل ولا
هرمون ولا استنساخ --- مثل \omterm{prop:b3:rna-regulation:translation}{المفتاح الريبوزي}، مستقلَب يستشعره الجزيء
الذي يتحكم فيه.
\textbf{17.} يُترجم الفيريتين بصرف النظر عن الحديد: ففيريتين المصل
عالٍ جدًا، وحديد المصل والمخزون سويّان (بلا فرط تحميل). وفي العدسة،
التي لا دوران فيها، يتراكم الفيريتين سنين ويتبلور --- فيكون ساد
باكر.
\textbf{18.} الدبوس موضع هبوط لا أكثر؛ والأثر هو ما يعترضه البروتين
المرتبط فيزيائيًا --- الريبوزومات الماسحة عند النهاية
5$'$، ونوكليازًا عند النهاية 3$'$.
\textbf{19.} كيمياءان تغطيان حالتين: فعنقود الحديد والكبريت في
IRP1 تهدمه المؤكسدات وقلة الأكسجين أيضًا، فيستجيب IRP1 لحالة
الأكسدة والإرجاع، في حين يستجيب تحلل IRP2 المعتمد على الحديد
للحديد نفسه على امتداد مدى الأكسجين الفيزيولوجي؛ والتكرار يجعل
الحُصَين متينًا، والمستشعران يبلّغان عن أمرين مختلفين.
\textbf{20.} $12\times 48\times 33\times 2 = \num{38016}$. و$1 - (1 -
50/\num{38016})^{50} = 1 - 0.99868^{50} \approx 1 - e^{-0.066} =
0.064$: أي نحو \qty{6}{\%}.
\textbf{21.} الأشكال المتطابقة ترتبط ارتباطًا مِثليًا فتتنافر: فتتجنب
تفرعات العصبون الواحد بعضها بعضًا فتنتشر، بينما قد تتقاطع تفرعات
عصبونين مختلفين لا يشتركان في شكل تقريبًا --- أي هوية خاصة لكل خلية.
\textbf{22.} $2^{n}$؛ و$\num{1024}$ عند $n = 10$. والعدد أقل في
الواقع: فالإدراج تقرره عوامل خاصة بالنمط الخلوي وهو منسّق
عبر الإكسونات، وكثير من التوليفات يزيح إطار القراءة فيهدمه
التحلل الموجَّه بالهُراء، والإكسونات ليست مستقلة.
\textbf{23.} تتحكم \emph{Sxl} أيضًا في \omterm{def:b3:chromatin-epigenetics:x-inactivation}{معاوضة الجرعة}: ففي الإناث
يحجب SXL ترجمة \emph{msl-2}، فلا يُفرط في استنساخ صبغيي X؛ ومن دون
SXL يركّب الجنين XX معقّد MSL على صبغيي X كليهما فيضاعف ناتجهما ---
أي جرعة قاتلة من كل مورّثة مرتبطة بالصبغي X، قبل أن يُطرح أمر الجنس أصلًا.
\textbf{24.} لا يعمل المروّج الباكر إلا ما دامت إشارة X:A حاضرة؛
وSXL الذي يصنعه يفرض التشذيب المنتج للنسخة الآتية من مروّج
الصيانة، وهو عامل في الجنسين، فمتى وُجد SXL ظل يصنع نفسه ولم تعد
الإشارة لازمة --- وهو منطق ذاتي الإدامة، \omterm{def:b3:chromatin-epigenetics:histone-code}{قارئ} يستدعي كاتبًا، كما
في علامات \omterm{def:b3:chromatin-epigenetics:nucleosome}{الصبغين}، لكن ببروتين وتشذيب هنا.
\textbf{25.} بروتين LIN-14 مكبوت $8$ أضعاف؛ وزمن نصف مفتاح الرسالة
\qty{10}{min}؛ وRNA غير المضخَّم يهبط دون جزيء واحد في الخلية بعد
$12$ انقسامًا.
\end{solution}
